% Complete Section 3.3, original PDF 33-37 / printed 23-27.
% block: adjoint-heading page:33 type:section
\subsection{{无界算子的伴随}}\label{sec:unbounded-adjoint}

% block: adjoint-functional-intro page:33 type:ordinary-paragraph
{现在设 $T$ 是从 $X$ 到 $Y$ 的任意稠密定义线性算子. 则可以定义一个从 $Y$ 到 $X$ 的无界线性算子, 称为 $T$ 的\emph{伴随}, 记为 $T^*$. 为此, 对任意 $w\in Y$, 考虑由下式给出的线性映射 $D(T)\to\mathbb R$:}

% block: adjoint-functional-equation page:33 type:equation
\begin{equation}\label{eq:adjoint-defining-functional}
v\mapsto\langle w,Tv\rangle_Y,\qquad v\in D(T).
\end{equation}

% block: adjoint-domain-definition page:33 type:ordinary-paragraph
{伴随的定义域 $D(T^*)$ 定义为使线性泛函 \eqref{eq:adjoint-defining-functional} 关于 $X$ 范数有界的 $w\in Y$ 的集合, 即存在常数 $c_w$ 使得}
\[
|\langle w,Tv\rangle_Y|\leq c_w\|v\|_X,\qquad v\in D(T).
\]

% block: adjoint-riesz-representation page:33 type:ordinary-paragraph
{对任意这样的 $w\in D(T^*)$, 映射 \eqref{eq:adjoint-defining-functional} 都唯一地延拓为从整个 $X$ 到 $\mathbb R$ 的有界线性泛函. 此外, 由 Riesz 表示定理, 存在唯一的元素 $x\in X$, 使得}
\begin{equation}\label{eq:adjoint-riesz-identity}
\langle w,Tv\rangle_Y=\langle x,v\rangle_X,\qquad v\in D(T),
\end{equation}

% block: adjoint-identity page:33 type:ordinary-paragraph
{并定义 $T^*w=x$. 简言之,}
\[
\langle T^*w,v\rangle_X=\langle w,Tv\rangle_Y,\qquad v\in D(T),
\]
{而 $w\in D(T^*)$ 当且仅当存在一个 $X$ 中的元素 (它必然唯一), 取其为 $T^*w$ 时能够满足此等式.}

% block: rotated-graph-definition page:34 type:ordinary-paragraph
{伴随的许多性质可由 $T$ 的图与其伴随 $T^*$ 的图之间的一个简单关系推出. 由于后者位于 $Y\times X$ 而非 $X\times Y$ 中, 我们将 $T^*$ 的图的一个变体, 即\emph{旋转图}, 定义为}
\[
\widetilde\Gamma(T^*)=\{\,(-x,y)\mid(y,x)\in\Gamma(T^*)\,\}=\{\,(-T^*w,w)\mid w\in D(T^*)\,\}\subset X\times Y,
\]
{其中加入负号是为了使这种关系更加明显. 事实上, $T^*$ 的旋转图就是 $\Gamma(T)^\perp$, 即图 $\Gamma(T)$ 在 $X\times Y$ 中的正交补.}

% block: adjoint-graph-proposition page:34 type:proposition
\Needspace{6\baselineskip}
\begin{Proposition}\label{prop:adjoint-rotated-graph}
{设 $T$ 为稠密定义的无界算子. 则}
\begin{equation}\label{eq:adjoint-rotated-graph}
\Gamma(T)^\perp=\widetilde\Gamma(T^*),
\end{equation}
\begin{equation}\label{eq:graph-closure-rotated-complement}
\overline{\Gamma(T)}=\widetilde\Gamma(T^*)^\perp.
\end{equation}
\end{Proposition}

% block: adjoint-graph-proof page:34 type:proof
\begin{Proof}
{$X\times Y$ 中的元素 $(x,w)$ 属于 $\Gamma(T)^\perp$, 当且仅当}
\[
\langle x,v\rangle_X+\langle w,Tv\rangle_Y=0,\qquad v\in D(T),
\]
{这正是 $w$ 属于 $D(T^*)$ 且 $T^*w=-x$ 的条件 \eqref{eq:adjoint-riesz-identity}, 也就是 $(x,w)$ 属于 $\widetilde\Gamma(T^*)$ 的条件. 这就证明了 \eqref{eq:adjoint-rotated-graph}. 再取正交补即得 \eqref{eq:graph-closure-rotated-complement}.}
\end{Proof}

% block: adjoint-closed-consequence page:34 type:ordinary-paragraph
{作为一个直接应用, 可知稠密定义无界算子 $T$ 的伴随 $T^*$ 总是闭的, 无论 $T$ 是否闭. 此外, 原算子 $T$ 是闭的当且仅当 $\Gamma(T)=\widetilde\Gamma(T^*)^\perp$, 因为 $\Gamma(T)=\overline{\Gamma(T)}$ 恰在 $T$ 闭时成立.}

% block: adjoint-dense-intro page:34 type:ordinary-paragraph
{作为 Proposition~\ref{prop:adjoint-rotated-graph} 的另一个应用, 我们证明\emph{闭的}稠密定义算子的伴随本身也稠密定义.}

% block: adjoint-densely-defined page:34 type:proposition
\Needspace{6\baselineskip}
\begin{Proposition}\label{prop:closed-adjoint-densely-defined}
{若 $T$ 是从 $X$ 到 $Y$ 的闭的稠密定义无界算子, 则 $T^*$ 是从 $Y$ 到 $X$ 的闭的稠密定义算子.}
\end{Proposition}

% block: adjoint-dense-proof page:34 type:proof
\begin{Proof}
{要证明 $T^*$ 稠密定义, 须证明任何与 $D(T^*)$ 正交的 $y\in Y$ 必为零. 现在有}
\[
y\perp D(T^*)\iff(0,y)\in\widetilde\Gamma(T^*)^\perp=\Gamma(T),
\]
{最后一个等式来自 Proposition~\ref{prop:adjoint-rotated-graph} 以及 $T$ 闭的假设. 因此, 若 $y\perp D(T^*)$, 则 $y=T0$, 因而确实为零.}
\end{Proof}

% block: double-adjoint-intro page:34 type:ordinary-paragraph
{由此可知, 当 $T$ 是闭的稠密定义算子时, $T^*$ 也是如此, 因而可以定义 $T^{**}$, 它仍是闭的稠密定义算子 $X\to Y$. 事实上, $T^{**}$ 与 $T$ 相同.}

% block: double-adjoint-proposition page:34 type:proposition
\Needspace{6\baselineskip}
\begin{Proposition}\label{prop:double-adjoint-identity}
{若 $T$ 是从 $X$ 到 $Y$ 的闭的稠密定义无界算子, 则 $T^{**}=T$.}
\end{Proposition}

% block: double-adjoint-proof page:34 type:proof
\begin{Proof}
{由于 $T$ 是闭的, 有 $\Gamma(T)=\widetilde\Gamma(T^*)^\perp$, 类似地, $\widetilde\Gamma(T^*)=\Gamma(T^{**})^\perp$. 将两式结合得到 $\Gamma(T)=\Gamma(T^{**})^{\perp\perp}$; 因为 $T^{**}$ 是闭的, 后者等于 $\Gamma(T^{**})$. $T$ 和 $T^{**}$ 的图相同, 因而两算子相等.}
\end{Proof}

% block: null-range-closed-intro page:35 type:ordinary-paragraph
{若 $T$ 是闭算子, 则其零空间 $\mathcal N(T)$ 是 $X$ 的闭子空间. 这直接由闭算子的定义得到. 然而, 闭线性算子的像空间完全可能不闭, 即使是有界线性算子也是如此 (例如, 考虑 $H^1$ 到 $L^2$ 的包含映射). 下面的定理概括了 $T$ 和 $T^*$ 的零空间与像空间之间的关系. 其证明是 Proposition~\ref{prop:adjoint-rotated-graph} 的又一个应用.}

% block: null-range-orthogonality page:35 type:theorem
\Needspace{6\baselineskip}
\begin{Theorem}\label{thm:adjoint-kernel-range-orthogonality}
{设 $T$ 是从 $X$ 到 $Y$ 的闭的稠密定义算子. 则}
\[
\begin{gathered}
\mathcal R(T)^\perp=\mathcal N(T^*),\quad\mathcal N(T)^\perp=\overline{\mathcal R(T^*)},\\
\mathcal R(T^*)^\perp=\mathcal N(T),\quad\mathcal N(T^*)^\perp=\overline{\mathcal R(T)}.
\end{gathered}
\]
\end{Theorem}

% block: null-range-orthogonality-proof page:35 type:proof
\begin{Proof}
{先验证第三个等式:}
\[
\begin{aligned}
z\in\mathcal N(T)&\iff(z,0)\in\Gamma(T)=\widetilde\Gamma(T^*)^\perp\\
&\iff\langle z,T^*w\rangle_X=0\ \forall w\in D(T^*)\iff z\in\mathcal R(T^*)^\perp.
\end{aligned}
\]
{将 $T$ 换成 $T^*$, 第三个等式便推出第一个等式. 分别对第一个和第三个等式取正交补, 即得第四个和第二个等式.}
\end{Proof}

% block: closed-range-theorem-intro page:35 type:ordinary-paragraph
{虽然闭线性算子 $T$ 的像空间可能不闭, 但若它是闭的, 则 $T^*$ 的像空间也闭 (反之亦然, 因为可以将 $T$ 换成 $T^*$ 应用这一结论). 这是 Banach 闭值域定理的推论. 在 Banach 空间上闭算子的更一般背景下, 证明可参见例如 [\BookNumber{71}, Chapter VII.5]. 为方便读者, 这里给出 Hilbert 空间情形的证明. 先陈述一个简单引理.}

% block: lower-bound-closed-range page:35 type:lemma
\Needspace{6\baselineskip}
\begin{Lemma}\label{lem:operator-lower-bound-closed-range}
{设 $T$ 是 Hilbert 空间 $X$ 和 $Y$ 之间的闭算子, 并假设存在 $\gamma>0$, 使得}
\[
\|Tv\|_Y\geq\gamma\|v\|_X,\qquad v\in D(T).
\]
{则 $T$ 的像空间是闭的.}
\end{Lemma}

% block: lower-bound-closed-range-proof page:35 type:proof
\begin{Proof}
{若 $\{v_n\}$ 是 $D(T)$ 中的一个序列, 且 $Tv_n$ 在 $Y$ 中收敛到某个 $y$, 须证明 $y\in\mathcal R(T)$. 假设意味着序列 $\{v_n\}$ 为 Cauchy 序列, 因而在 $X$ 中收敛到某个 $x$. 由于 $T$ 闭, 有 $x\in D(T)$ 且 $y=Tx\in\mathcal R(T)$.}
\end{Proof}

% block: closed-range-theorem page:35 type:theorem
\Needspace{6\baselineskip}
\begin{Theorem}{{闭值域定理}}\label{thm:closed-range}
{若从 $X$ 到 $Y$ 的闭的稠密定义算子 $T$ 的像空间在 $Y$ 中闭, 则 $T^*$ 的像空间在 $X$ 中闭.}
\end{Theorem}

% block: closed-range-proof-restriction pages:35,36 type:proof
\begin{Proof}
{设 $W=\mathcal R(T)$, 它是 $Y$ 的闭子空间; 设 $\overline T$ 为 $T$ 的像空间限制, 即从 $X$ 到 $W$ 的无界算子, 它与 $T$ 的定义域相同, 并且对 $v\in D(T)$ 有 $\overline Tv=Tv$. 显然, $\overline T$ 也是 Hilbert 空间之间的闭的稠密定义算子, 而且为满射. 我们将证明 $\mathcal R(\overline T^*)=\mathcal R(T^*)$. 这由两点观察得出. 首先注意 $D(\overline T^*)=D(T^*)\cap W$, 且 $\overline T^*$ 是 $T^*$ 在此空间上的限制. 这直接由伴随算子的定义以及对 $D(T)=D(\overline T)$ 中的所有 $v$ 有 $\overline Tv=Tv$ 得到. 因此 $\mathcal R(\overline T^*)\subset\mathcal R(T^*)$. 但 $W^\perp=\mathcal N(T^*)$. 所以, 若 $y\in D(T^*)$, 且 $w$ 为其在 $W$ 上的正交投影, 则 $w\in D(\overline T^*)$ 且 $\overline T^*w=T^*y$, 这表明 $\mathcal R(T^*)\subset\mathcal R(\overline T^*)$.}

% block: closed-range-proof-surjection page:36 type:proof
{由此, 只须证明: 若 $T$ 是闭的稠密定义\emph{满射}, 则 $\mathcal R(T^*)$ 是闭的. 我们知道, $D(T)$ 配备图范数后是 Hilbert 空间. 设 $Z$ 为 $\mathcal N(T)$ 在 $D(T)$ 中的正交补, 并设 $\widetilde T:Z\to Y$ 为 $T$ 的限制. 则 $\widetilde T$ 是一一且到上的有界线性算子. 因而它可逆, 所以存在正常数 $c$, 使对任意 $y\in Y$, 存在唯一 $z\in Z$ 满足 $Tz=y$ 且 $\|z\|_X\leq c\|y\|_Y$. 对 $y\in D(T^*)$ 应用这一点, 有}
\[
\|y\|_Y^2=\langle y,Tz\rangle_Y=\langle T^*y,z\rangle_X\leq\|T^*y\|_X\|z\|_X\leq c\|y\|_Y\|T^*y\|_X.
\]
{因此}
\[
\|T^*y\|_X\geq c^{-1}\|y\|_Y,\qquad y\in D(T^*).
\]
{现将前面的引理应用于 $T^*$, 即得 $\mathcal R(T^*)$ 是闭的.}
\end{Proof}

% block: finite-codimension-intro page:36 type:ordinary-paragraph
{本节最后给出一些验证算子像空间为闭的有用结论. 首先证明, 若闭线性算子的像空间为有限余维, 则它是闭的. 有界算子情形的证明见 [\BookNumber{52}, Lemma 19.1.1], 且很容易推广到闭算子.}

% block: finite-codimension-closed-range page:36 type:theorem
\Needspace{6\baselineskip}
\begin{Theorem}\label{thm:finite-codimension-closed-range}
{设 $T$ 为 Hilbert 空间 $X$ 和 $Y$ 之间的闭线性算子. 若商空间 $Y/\mathcal R(T)$ 为有限维, 则 $T$ 的像空间在 $Y$ 中闭.}
\end{Theorem}

% block: finite-codimension-proof page:36 type:proof
\begin{Proof}
{先考虑 $T$ 为从 $X$ 到 $Y$ 的有界线性算子的情形. 将 $T$ 限制在其零空间的正交补上, 可以假设 $T$ 为单射. 取 $y_1,\ldots,y_n\in Y$, 使它们模 $\mathcal R(T)$ 的陪集构成 $Y/\mathcal R(T)$ 的一组基. 则 $Y$ 的每个元素都可以唯一地写成 $r+c_1y_1+\cdots c_ny_n$, 其中 $r\in\mathcal R(T)$, $c_i\in\mathbb R$. 因此, 若通过下式定义 $\widetilde T:X\times\mathbb R^n\to Y$:}
\[
\widetilde T(x,c_1,\ldots,c_n)=Tx+c_1y_1+\cdots c_ny_n,
\]
{则 $\widetilde T$ 为一一且到上的有界线性算子, 从而由 Banach 定理可知它是同胚. 由于 $X\times\{0\}$ 在 $X\times\mathbb R^n$ 中闭, 可知 $\mathcal R(T)=\widetilde T(X\times\{0\})$ 在 $Y$ 中闭.}

% block: finite-codimension-proof-unbounded page:36 type:proof
{对从 $X$ 到 $Y$ 的闭线性算子 $T$ 的一般情形, 可将 $T$ 视为从 $D(T)$ 到 $Y$ 的有界线性算子, 因而由前一情形可知, 若其像空间为有限余维, 则它是闭的.}
\end{Proof}

% block: compactness-criterion-intro page:36 type:ordinary-paragraph
{第二个闭值域判据基于紧性.}

% block: compact-domain-closed-range page:36 type:theorem
\Needspace{6\baselineskip}
\begin{Theorem}\label{thm:compact-domain-closed-range}
{设 $T$ 为 Hilbert 空间 $X$ 和 $Y$ 之间的闭线性算子, 并假设其定义域到 $X$ 的包含映射是紧的 (定义域配备图范数). 则 $\mathcal R(T)$ 在 $Y$ 中闭.}
\end{Theorem}

% block: compact-domain-proof-boundedness pages:36,37 type:proof
\begin{Proof}
{给定元素序列 $x_n\in D(T)$ 和 $y\in Y$, 满足在 $Y$ 中有 $Tx_n\to y$, 须证明 $y\in\mathcal R(T)$. 不失一般性, 可取 $x_n\in\mathcal N(T)^\perp$. 此时序列 $x_n$ 在 $X$ 中有界, 下面用反证法说明. 事实上, 若非如此, 则可选取子序列, 仍记为 $x_n$, 使 $\|x_n\|_X\to\infty$. 定义 $w_n=x_n/\|x_n\|_X\in D(T)\cap\mathcal N(T)^\perp$, 则得到满足 $\|w_n\|_X=1$ 且 $Tw_n\to0$ 的序列. 这样的序列在 $D(T)$ 中有界, 故由紧性可选取子序列, 关于 $X$ 的范数收敛到某个 $w\in X$. 由于 $T$ 闭, 所得 $w$ 属于 $D(T)$ 且满足 $Tw=0$, 即属于 $\mathcal N(T)$. 然而, 作为 $w_n$ 的极限, $w$ 还属于 $\mathcal N(T)^\perp$ 且范数为 $1$. 这些条件显然矛盾.}

% block: compact-domain-proof-limit page:37 type:proof
{因此, 只须证明: 对元素的有界序列 $x_n\in D(T)$, 若在 $Y$ 中有 $\lim Tx_n=y$, 则 $y$ 属于 $\mathcal R(T)$. 再次利用紧性假设, 将该序列替换为一个子序列, 使 $x_n$ 还在 $X$ 中收敛到某个 $x$. 由于 $T$ 是闭的, 有 $x\in D(T)$ 且 $Tx=y$, 即得所需结论.}
\end{Proof}
